\documentclass[11pt,a4paper]{article}
\usepackage[T1]{fontenc}
\usepackage{lmodern,amsmath,amssymb,amsthm,booktabs,microtype,hyperref}
\usepackage[margin=28mm]{geometry}
\setlength{\emergencystretch}{3em}
\newtheorem{proposition}{Proposition}
\newtheorem{lemma}[proposition]{Lemma}
\newtheorem{theorem}[proposition]{Theorem}
\newcommand{\Q}{\mathbb Q}\newcommand{\Z}{\mathbb Z}
\title{A quartic in 193 rational coordinates}
\author{}\date{17 November 2026}

\begin{document}\maketitle

\begin{abstract}
Under the explicitly retained elliptic, contact-height, pole-realization
and simultaneous-prime premises, we construct an integer quartic in
193 total rational coordinates whose input projection is the integers
for every sufficiently large fixed parameter. The polynomial is the
sum of 167 primitive quadratic squares, has 2770 monomials, and has
exact coefficient height
$\max\{126616122944180358034867798936,149^{2N+2}\}$.
An equivalent reference basis has 2626 monomials. Polynomial cross
products, a cleared input numerator and a homogeneous Lucas circuit
remove stored ratios; an invertible Laurent transformation makes every
equation quadratic. We prove extraction at every odd contact depth,
integer completeness through synchronized formal depths, and coverage
of the permitted zero-ordinate and finite-torsion boundaries. A diagonal change of twelve coordinates reduces the parameter degree
of the quartic from four to two. Complete coefficient tables and
independent exact checks specify the families.
The sufficient defining parameter remains ineffective.
\end{abstract}

\paragraph{Recovery and provenance.}
This version reconstructs the 193-coordinate result from the published
194-coordinate system of 17 November 2026 (v73). The interrupted
working files were not available. All coefficient tables, Laurent
weights and verification receipts in this version are regenerated;
no byte-for-byte recovery of the interrupted tables is asserted.
The physical system and its arithmetic proof are unchanged. The new
steps are the seven-coordinate unit chart and the parameter rescaling
proved below.


\section{Arithmetic statement and retained premises}

Let $\mathcal S_N$ be the physical inverse-Vieta source
of~\cite{InverseVieta}, with its unchanged scalar, reference, contact,
trace and pole equations, and with the point charts, comparisons,
input test, Lucas circuit and factor assignments specified below.
The parameter is $149^N$, with $N\ge1$ fixed before the rational input.
The retained arithmetic premises are the following exact source
interfaces:
\begin{enumerate}
\item The group description for
$K=\Q(\sqrt2)$, $d=3-\sqrt2$, $D=d^2$, and
$E:y^2=x^3-Dx$, with $P_0=2R_0$,
$R_0=(1+\sqrt2,2\sqrt2)$:
$E(K)=\Z R_0+E[2](K)$, with unique free coefficient and all rational
torsion of order dividing two. The source reference chart, its
nonvanishing conclusions, and its formal logarithm with relative
error of order four are retained.
\item The scalar extraction and uniform ordered scalar preparation of
\cite{InverseVieta}, using the interfaces P04 and P11--P14 in
\cite{Uniform}. With $a=2\nu(Q_2)$ and $c=|a|$, they give
\begin{equation}\label{eq:scalar-bound}
 \begin{gathered}
 c\ge149^{L_N},\qquad
 h(\beta)\le 2^{330770}149^{-L_N}c(1+\log c/3),\\
 L_N=\left\lfloor\frac{9(2N-1)}4\right\rfloor+1,
 \qquad \beta=51X^2+2504.
 \end{gathered}
\end{equation}
\item The contact-height interface for
$\Theta=\Theta_0\cup(2-\Theta_0)$,
$\Theta_0=\{0,-1/2,-1,1/2399,-32/407\}$, and
$\Sigma=\{p:p\le76176\}\cup\{2029147\}$.
For $t\notin\Theta$, let $M(t)$ be the squarefree product of the
primes outside $\Sigma$ for which some $v_p(t-\theta)>0$ is odd.
The retained assertion is $h(t)\le H_2 M(t)^{7/6}$ with fixed $H_2$,
including $t=1$ and $M(1)=1$; see~\cite{Quartic204,Uniform}.
\item The positive weighted trace-height theorem of~\cite{Weighted},
used with all weights one, on the full finite-point domain including
finite torsion translates. It is used in exactly the form stated in
\eqref{eq:trace-height} below.
\item The pole obstruction and prepared actual-input realization
retained in~\cite{InverseVieta,Poles,Uniform}. The obstruction forbids
a negative odd valuation outside $\Sigma$. The prepared realization
applies to each actual pole input after its positivity, upper bound,
square reduced denominator and the source local conditions at
$2,3,17$ have been established. Its analytic and arithmetic foundations
are included in this retained interface.
\item The fixed-data simultaneous Gaussian prime-pattern supplier,
its compatible annihilation congruences, and its ordered real and
finite-place preparation, as used in~\cite{InverseVieta,Uniform}.
The prime supplier allows arbitrarily large prescribed slope and
prime lower bounds for each fixed finite coefficient list. No
uniformity in the input-dependent list is assumed.
The all-order separation is the consequence already established in
\cite[Theorem 2.2]{PrimeSeparation} and used in~\cite[P16]{Uniform}.
\end{enumerate}

\begin{theorem}\label{thm:arithmetic}
Under these retained premises there is a fixed integer $N_0$ such
that, for every $N\ge N_0$, the rational input projection of
$\mathcal S_N$ is exactly $\Z$. The sufficient $N_0$ is ineffective.
The Laurent compilation of this physical system inherits the same
arithmetic assertion whenever its exact row identities, nonzero unit
boundary and input-preserving inverse map are certified.
\end{theorem}

The compiled presentation uses the 191-coordinate physical storage
with $A=L_2$ retained; the final unit chart replaces five source unit
coordinates by seven, a net addition of two. Its
193-coordinate census and quadratic coefficient table are algebraic
certificates separate from the arithmetic proof. No deletion of $A$
is part of the system considered here.

The unchanged scalar equations first give
\[
 r_s=\beta(\beta-1)(\beta-2401)(32\beta+343)>0,\qquad
 R=\frac{149^{2N+1}r_s}{1+149^2r_s}>100.
\]
The forced-prime argument then excludes the reference torsion and
indices $-1,0,1$ before any elliptic logarithm or ordinate ratio is
used. These statements use only the retained scalar/reference rows.

\section{A finite-pair chart and its boundaries}

Use the following chart for the point pairs with roles $1,3,4$;
role $2$ keeps the inverse-Vieta scalar chart.
Work over a field of characteristic zero on $E:y^2=x^3-Dx$, $D\ne0$.
Take $s,t,F,G,h,d$ and put $p=h+s+t$. Consider
\begin{equation}\label{eq:chart}
 d=pF=hG,\qquad D(F+G)=(s-t)h-t^2,\qquad
                       1=d(s-t)-DFG.
\end{equation}

\begin{proposition}
Every solution has $FG\ne0$ and gives distinct finite points
\[
                 U=(F^{-1},s/F),\qquad V=(G^{-1},t/G)
\]
on $E$. Their difference is finite and has abscissa $hp$.
Conversely every pair of distinct finite points with nonzero abscissas
supplies the chart, including the tangent case $V=-U$.
Finite two-torsion differences are included; $d$ need not be nonzero.
\end{proposition}
\begin{proof}
For the four residuals
\[
 r_1=d-pF,\quad r_2=d-hG,\quad
 r_3=D(F+G)-(s-t)h+t^2,\quad r_4=d(s-t)-1-DFG,
\]
there are exact polynomial identities
\[
 Fs^2+DF^2-1=r_4-(s-t)r_1+Fr_3,
\quad
 Gt^2+DG^2-1=r_4-(s-t)r_2+Gr_3.
\]
They give $F,G\ne0$ and both curve equations.
If $F\ne G$, the chord joining $U$ to $-V$ has slope $s+h$:
$h(G-F)=F(s+t)$. The chord formula gives
\[
 x(U-V)=(s+h)^2-F^{-1}-G^{-1}=h(h+s+t)=hp.
\]
If $F=G$, the equations give $t=-s$ and
$2sh=F^{-1}+DF$. One has $s\ne0$, and $s+h$ is the tangent slope
at $U$. Hence the same conclusion follows by the tangent law.
In particular $U=V$ is impossible: it would force $s=t=0$ in this
last case, contrary to the displayed equation and the curve equation.
Reversing the chord or tangent calculation proves the converse.
If $d=0$, then $h=p=0$ and $x(U-V)=0$; this is the allowed finite
two-torsion difference, rather than an identity difference.
\end{proof}

The new chart removes numerator identity points. The retained
completeness construction uses nonzero nontorsion numerator multiples
and can avoid their finitely many forbidden equalities, so that
restriction causes no loss there. The retained second primary ratio
and scalar conclusion $S\ne0$ force $s_2\ne0$. The final polynomial
input and Lucas tests allow zero ordinates in the other numerator roles.
They are addressed by local extraction and the global height argument,
not by unstated inverse guards.

\section{The physical comparisons}

Let $K=\Q(\sqrt2)$ and use the source fixed point $P_0=2R_0$.
The unchanged reference chart gives
\[
 f_H=1/x(H)\ne0,\qquad S=y(H)/x(H)\ne0,\qquad
 T_H=y(H-P_0)/x(H-P_0)\ne0.
\]
For the numerator pair $U_i,V_i$, put
$s_i=y(U_i)/x(U_i)$ and $t_i=y(V_i)/x(V_i)$.
Let $(k_i,c_i)=(1,0),(2,0),(2,0),(4,1)$ for $i=1,2,3,4$.
At places where the indicated denominators are nonzero, define the
primary ratios
\[
 L_i=k_iS/s_i+c_i.
\]
Define the polynomial cross products
\[
                         E_i=S\overline{t_i}-\overline{T_H}s_i.
\]
Whenever $s_i t_i\ne0$, the former comparison is exactly
\begin{equation}\label{eq:cross}
 L_i-\overline{\left(k_iT_H/t_i+c_i\right)}
                      =\frac{k_iE_i}{s_i\overline{t_i}}.
\end{equation}
Denote the left side of~\eqref{eq:cross} by $C_i$ whenever its
ratios are defined.
The affine constant cancels. For the unchanged scalar inverse-Vieta
chart, $s_2=(U+M)/R$ and $t_2=-(V+M)/R$, and the corresponding
quadratic equation is
\[
 R E_2=-S\overline{(V+M)}-\overline{T_H}(U+M),
\]
with the retained rational $R\ne0$.

Keep $A=L_2$ and the scalar primary equation
$A(U+M)=2RS$. The input ratio coordinate is replaced by the polynomial test
\[
                         \mathcal N=S-\beta s_1.
\]
At a place with $s_1\ne0$, the original input test is recovered as
$I=L_1-\beta=\mathcal N/s_1$.

Put $w=s_4$. The Lucas circuit uses four field coordinates
$C,P,Q,\mathcal H$ and the quadratic equations
\begin{equation}\label{eq:homogeneous}
 \begin{aligned}
 C&=(A-1)(w+4S),& P s_3&=2S w,\\
 Q&=(A+1)P,& \mathcal H&=P^2+C^2-CQ-w^2.
 \end{aligned}
\end{equation}
At a place where $s_3w\ne0$, define
$L_3=P/w=2S/s_3$ and $L_4=1+4S/w$.
For the original Lucas polynomial
\[
 \mathcal G=L_3^2-(L_2+1)(L_2-1)L_3L_4+(L_2-1)^2L_4^2-1.
\]
there is the exact identity
\begin{equation}\label{eq:homogeneous-identity}
                         \mathcal H=w^2\mathcal G.
\end{equation}
No separate $L_1,L_3,L_4$ or quotient $\mathcal G/f_H$ coordinate
is required. The displayed equations recover the original primary
ratios at every place at which they are used in the extraction argument.
There is no global division by $s_1,s_3$ or $w$ in the new equations.

\section{The five factors for each detector}

Write $D_H=f(H)$, $B_0=\overline{f(H-P_0)}$,
$A_i=f(U_i)$ and $B_i=\overline{f(V_i)}$, where $f(Q)=1/x(Q)$.
The conjugations in the $B$ tests are part of the physical definition.
The two contact arguments are $t_0=\lambda$ and $t_1=\lambda+T$.
Their physical weights are $q_0=u$ and $q_1=uv$, respectively.
Each detector has four ordinary factors
\[
 F_{j,k}=1+q_j a_{j,k}+q_j^2 b_{j,k}+q_j^3 c_{j,k}
 \qquad(0\le k\le3)
\]
and one special factor
\[
 F_{j,4}=1+q_j^4\mathcal N^8
                +q_j a_{j,4}+q_j^2 b_{j,4}+q_j^3 c_{j,4}.
\]
The exact assignment is
\begin{center}
\begin{tabular}{ccccc}
\toprule
$j$&$k$&$a_{j,k}$&$b_{j,k}$&$c_{j,k}$\\
\midrule
0&0&$B_3^2$&$D_H^2$&$E_2^4$\\
0&1&$A_1^2$&$B_0^2$&$B_1^2$\\
0&2&$A_3^2$&$B_4^2$&$E_3^4$\\
0&3&$\mathcal H^4$&$B_2^2$&$E_4^4$\\
0&4&$A_4^2$&$A_2^2$&$E_1^4$\\
\midrule
1&0&$E_2^4$&$E_1^4$&$E_4^4$\\
1&1&$B_3^2$&$D_H^2$&$A_3^2$\\
1&2&$B_1^2$&$B_4^2$&$E_3^4$\\
1&3&$A_1^2$&$B_2^2$&$\mathcal H^4$\\
1&4&$B_0^2$&$A_4^2$&$A_2^2$\\
\bottomrule
\end{tabular}
\end{center}
Thus each detector tests each of the ten reciprocal abscissas and
each of $E_1,E_2,E_3,E_4,\mathcal H,\mathcal N$ once. Put
$W_j=\prod_{k=0}^4F_{j,k}$ and use the unchanged norm, reciprocal-trace
and pole equations. Whenever the denominators are nonzero, the actual
trace is $V_j=\operatorname{Tr}_{K/\Q}(W_j^{-1})$ and the actual
pole input is $\gamma^2 R(t_j)/V_j^2$, where $\gamma=408/409$
and $q_j=R(t_j)^{-1}-1$ are the source conventions.

\section{Extraction at arbitrary negative odd weight depth}

Let $p>5$ be a good unramified consumed prime, and suppose
$v_p(q)=-\mu<0$ with $\mu$ odd. Use the two conjugate valuation
slots, which may coincide when $p$ is inert in $K$. In either slot, every
finite reciprocal-abscissa valuation is even. This follows either
from $h^2=f-Df^3$ or from $fs^2+Df^2=1$.

Every coefficient $a_{j,k},b_{j,k},c_{j,k}$ is a fourth power or
the square of a reciprocal abscissa, so its finite valuation is a
multiple of four. Zero tests are allowed and are omitted when taking
the minimum of term valuations.

\begin{lemma}
Every factor is nonzero, has valuation at most zero, and has inverse
residue zero or one. Under the unchanged trace and pole equations,
there is one valuation slot above $p$ where every reciprocal test is positive,
every $E_i$ and $\mathcal H$ is positive, and $v(\mathcal N)\ge1$.
\end{lemma}
\begin{proof}
For an ordinary factor the valuation classes modulo four are
\[
                     0,-\mu,-2\mu,-3\mu.
\]
They are distinct because $\mu$ is odd. For the special factor,
the disjoint classes modulo eight are
\[
 \{0\},\qquad\{4\},\qquad
 \{-e\mu,-e\mu+4\}\quad(e=1,2,3).
\]
These are the classes of the retained five-factor lemma.
Hence there is a unique minimum.
It is at most zero, and a zero minimum has constant term one and
all other terms of positive valuation.

The contact identities give $v_p(R(t_j))=\mu$. In both conjugate slots the
five factors are nonzero by their unique minimum, so $W_j\ne0$.
The unchanged norm-over-trace equation then forces $V_j\ne0$.
If $v_p(V_j)\ge\mu/2$, its integral valuation and the oddness of
$\mu$ would give the actual pole input negative odd valuation,
contrary to the retained pole obstruction. Hence $v_p(V_j)<\mu/2$.
At least one of the two summands of $V_j$ has valuation less than
$\mu/2$, so in one slot $v(W_j)>-\mu/2$. Each factor has valuation
at most zero, so every factor at that place has valuation greater
than $-\mu/2$. Consequently a summand $q^eT^m$ satisfies
\[
                          v(T)>\frac{e-1/2}{m}\mu.
\]
For fourth-power tests and $e=1,2,3$, these coefficients are
$1/8,3/8,5/8$, all positive. For reciprocal squares they are
$1/4,3/4,5/4$; even valuation then gives depth at least two.
For $\mathcal N$ the coefficient is $7/16$, giving depth at least one.
\end{proof}

At that successful place, positive reciprocal tests give
\[
 v(s_i)\le-1,\qquad v(\overline{t_i})\le-1,
 \qquad v(f_H)\ge2.
\]
In particular every $s_i,t_i$ is nonzero and the local use of
\eqref{eq:cross} is legitimate. It gives comparison depth at least
three. Since $v(w)\le-1$, equation~\eqref{eq:homogeneous-identity}
gives $v(\mathcal G)=v(\mathcal H)-2v(w)\ge3$.
Similarly $v(I)=v(\mathcal N)-v(s_1)\ge2$.

The original integrality transfer must still precede absolute
formal-logarithm precision. At least one of the reference indices
$h,h-1$ is a $p$-adic unit. The corresponding exact logarithmic
ratio is integral; comparison depth three transfers integrality to
the other actual ratio. Its relative order-four error then also
becomes an absolute error of order at least four. Thus the four
primary ratios have the original label congruences modulo $p^3$,
and the input comparison has precision at least two. The original radical
divisibilities, integer Lucas residual and height soundness follow
without a lost contact or an altered exponent.

\begin{proposition}\label{prop:input-depth-two}
For each detector and every rational solution of the physical system,
if $J_X=|\operatorname{num}(\beta-n)|$ is nonzero, then
\[
                              M_j^2\mid J_X.
\]
\end{proposition}
\begin{proof}
At each prime $p\mid M_j$, the successful valuation slot has
$v(\mathcal N)\ge1$ and $v(s_1)\le-1$. Consequently
$v(I)=v(\mathcal N/s_1)\ge2$, whereas $L_1\equiv n\pmod{p^3}$.
It follows that $v_p(\beta-n)\ge2$. Both $n\in\tfrac12\Z$ and
$L_1$ are $p$-integral, so the reduced denominator introduces no
$p$-factor. Taking the product over the squarefree radical proves
the assertion. If $M_j=1$, it is automatic.
\end{proof}
The height proof below uses only the original weaker $M_j\mid J_X$.

If a detector has no consumed prime, its radical is one and these
divisibilities are automatic. No globally defined alternate ratio
is required merely to state that trivial divisibility. In particular
global zero-ordinate strata are not discarded without a reason.
The same observation covers any zero product or zero-trace boundary
left by a cleared detector equation: the calculation at a consumed
prime proves both the product and actual trace nonzero before dividing.
Consequently such a boundary can occur only in a detector whose
contact radical is one.
The scalar primary ratio remains globally defined by
$A(U+M)=2RS$, with $R,S\ne0$. There is no implicit division by
$P$ in this system. The only new denominator exclusions used locally
are $s_1,s_3,s_4\ne0$, each supplied by a reciprocal test.

The retained inverse rows force $q_0,q_1\ne0$. The contact equations
exclude every $t_j\in\Theta$ without division, but allow $t_j=1$,
where $q_j=-1$ and $M_j=1$. No division by $t_j-1$ is made in that
case. The full-domain contact-height premise includes this argument,
so it is covered by the same height proof.

For completeness of the global argument, all four difference points
$Q_i$ remain finite under the new chart, including possible finite
torsion differences. Hence the retained trace-height theorem still gives
\begin{equation}\label{eq:trace-height}
 \begin{gathered}
 B^2\le A_0h(T)+J_0,\qquad
 B=\max(1,|n|,|a|,|t|,|b|),\\
 T=7+\tfrac14\operatorname{Tr}\!\left(\sum_{i=1}^4x(Q_i)\right).
 \end{gathered}
\end{equation}
This estimate depends only on the points and their labels. It does not
use a globally stored $L_3$ or any of the implicit ratios. At every
consumed prime the local recovery above proves the original divisibilities
$M_j^3\mid\Delta$ and $M_j\mid J_X$ (when $J_X\ne0$). They are
automatic when $M_j=1$, so the original two-contact height argument
applies to all rational solutions. In particular if any $s_1,s_3,s_4$
or alternate denominator vanishes globally, no prime can be consumed
by either detector. The two radicals are then one; the unchanged
contact-height bounds make $h(T)$, and hence $B$, bounded. This
contradicts the retained scalar lower bound $|a|\ge149^{L_N}$ for
the same sufficiently large fixed $N$. No revised numerical or
asymptotic height threshold is asserted.

Deleting the alternate scalar ratio equation permits $M+V=0$,
and hence a two-torsion scalar numerator $V_2$. The retained primary
equation still forces $M+U\ne0$, and the scalar quotient equation
forces $M+W\ne0$. Thus the scalar root $Q_2$ remains nontorsion and
its scalar extraction is unchanged. The Vieta point identities also
hold when the alternate numerator has zero ordinate.

\section{The global height conclusion}

The group premise and the finite difference charts give the labels
\[
 n=\nu(Q_1)\in\tfrac12\Z,\qquad
 a=2\nu(Q_2)\in\Z\setminus\{0\},\qquad
 t=2\nu(Q_3)\in\Z,\qquad
 b=4\nu(Q_4)+1\in2\Z+1.
\]
In particular $b\ne0$ even when $Q_4$ is finite torsion.
At both real embeddings, every finite point satisfies $x\ge-d$
with the corresponding value of $d$. Therefore the displayed
aggregate has $T\ge1$. Put $M_j=M(t_j)$ and
\[
 \Delta=t^2-(a+1)(a-1)tb+(a-1)^2b^2-1\in\Z.
\]
The local extraction proves $M_j^3\mid\Delta$ and, when nonzero,
$M_j\mid J_X=|\operatorname{num}(\beta-n)|$. The bounds
$|\Delta|\le14B^4$ and $J_X\le4B e^{h(\beta)}$ are unchanged.
The two separate divisibilities do not require $M_0,M_1$ to be
coprime.

For any nonzero integer $J$ with $M_j^k\mid J$ for $j=0,1$, the
contact-height bound, $h(T)\le h(t_0)+h(t_1)+\log2$, and
\eqref{eq:trace-height} give fixed constants with
\[
                         B^2\le C|J|^{7/(6k)}+J_1.
\]
If $\Delta\ne0$, taking $k=3$ gives
$B^2\le C_1B^{14/9}+J_1$, so $B$ is bounded by a fixed constant.
The scalar lower bound~\eqref{eq:scalar-bound} excludes this case
for all sufficiently large fixed $N$.

When $\Delta=0$, the retained elementary Lucas argument, with
parameter $|a+1|$, gives
\[
                         \log B\ge(c-6)\log(c-3)
\]
for sufficiently large $c=|a|$. This includes the allowed zero
third label: if $t=0$, the equation would force
$(a-1)^2b^2=1$, which is incompatible with large $|a|$.
If $J_X\ne0$, the preceding height estimate with the original
choice $k=1$ gives
\[
                 \log B\le \tfrac75h(\beta)+C_2.
\]
Together with~\eqref{eq:scalar-bound}, this contradicts the Lucas
lower bound for every sufficiently large fixed $N$. Thus
$\beta=n\in\tfrac12\Z$. If $X=r/s$ is reduced, this forces
$s^2\mid102$; since $102$ is squarefree, $s=1$. This proves soundness.
Only the original input precision one is used in this height argument;
the stronger precision two does not change its threshold calculation.

\section{Fresh integer witnesses and synchronized formal depths}

Fix an integer $X$ and a parameter beyond the uniform scalar
completeness threshold in~\cite{InverseVieta}. The retained ordered
scalar construction first chooses its finite and real preparation
and then an arbitrarily large positive $a\in6\Z$. Its actual scalar
pole input satisfies the source prepared-realization conditions.
The factor $3$ in this extra divisibility is fixed in advance and is
permitted by the retained compact-group progression argument.

For the Lucas recurrence $U_0=0$, $U_1=1$,
$U_{r+1}=(a+1)U_r-U_{r-1}$, put
\[
                   b=\frac{U_{a-1}}{a-1},\qquad t=U_a.
\]
Modulo $a-1$ one has $U_r\equiv r$, so $b$ is integral. Since
$a+1$ is odd, the recurrence modulo two has period $0,1,1$.
Here $a-1\equiv2\pmod3$ and $a-1$ is odd; hence $b$ is odd.
The adjacent-term Lucas identity gives $\Delta=0$, and taking $a$
large gives $b>1$. Fix
\[
 Q_1=[\beta]P_0,\qquad Q_2=[a]R_0,\qquad
 Q_3=[t]R_0,\qquad Q_4=[(b-1)/2]R_0.
\]
All four are nonzero nontorsion points. Their aggregate $T$, all
root multipliers, and the source finite and real neighborhoods are
fixed before comparison primes are chosen. For an integer reference
index $h$, take $H=[h]P_0$, $U_i=[h]Q_i$ and $V_i=[h-1]Q_i$.
The point differences are exactly $Q_i$, so the new charts supply
rational coordinates whenever $h\ne0,1$ and the source neighborhoods
avoid their finitely many forbidden finite points. Their ratio labels
are exactly $\beta,a,t,b$ in the formal logarithms.

Unequal reference depths cannot be ignored. If the primary depth is
$d$ and the opposite depth is $e$, the ordinary formal estimates give
only
\[
                        v(E_i)\ge4\min(d,e)-d-e,
\]
when the fixed numerator multipliers are local units. This can be
negative. The required repair is a refinement of the reference
progression after the comparison primes are known.

\begin{lemma}\label{lem:sync}
The retained simultaneous-prime input permits the comparison primes
to be chosen so that none divides any selected elliptic reduction-group
order or the previously fixed preparation modulus. At each of those
primes, the reference progression can then be refined so that the
primary and opposite reference formal depths are equal, while
retaining every original annihilation congruence and real return.
\end{lemma}
\begin{proof}
The separation assertion is supplied by
\cite[Theorem 2.2]{PrimeSeparation}, under precisely the retained
weighted prime-pattern input, and by its application in
\cite[P16]{Uniform}. In this common-leading additive family its
needed numerical consequence also follows directly from a short
interval and full two-torsion, as follows.

The fixed contact lists at $\lambda$ and $\lambda+T$, with $T$
fixed before selection, become one finite family
$p_i=Z+c_iL^2$ after clearing rational denominators. Repeated
coefficients are merged with their incidences retained. The retained
prime theorem permits an arbitrarily large prescribed lower bound
on $Z/L^2$, so all ratios $p_i/p_j$ may be made arbitrarily close
to one. Here $Z$ is the prime-pattern parameter and is unrelated to
the fixed exponent $N$. At both places of each split selected prime the reduction
has its four two-torsion points; its order is divisible by four.
Hasse's bound gives
\[
                 \max_j\#E(\mathbb F_{p_j})<2\min_i p_i
\]
for sufficiently large selected patterns, in either conjugate model.
If $p_i$ divided one such order, that order would equal the odd
prime $p_i$, contradicting divisibility by four. The cutoff also
avoids every fixed preparation modulus and every numerator or
denominator of a fixed root multiplier.

For a selected split $p$, let $v_+,v_-$ be the paired valuations
normalized by $v_+(p)=v_-(p)=1$. Let $o_+,o_-$ be the actual
reduction orders of $P_0$ at those completions. They are prime to
every selected prime.
The retained congruences make $h$ divisible by $o_+$ and $h-1$
divisible by $o_-$. Put
\[
 d_0=v_+\bigl(\zeta([o_+]P_0)\bigr),\qquad
 e_0=v_-\bigl(\zeta([o_-]P_0)\bigr),\qquad \zeta=-x/y.
\]
These are positive finite integers. The formal logarithm is an
isometry on the good formal subgroup at $p>5$, so the primary and
opposite depths are
\[
                    d_0+v_p(h),\qquad e_0+v_p(h-1).
\]
If $d_0<e_0$, impose $v_p(h)=e_0-d_0$; then $h-1$ is a unit.
If $e_0<d_0$, impose $v_p(h-1)=d_0-e_0$; then $h$ is a unit.
If the two are equal, require $h\not\equiv0,1\pmod p$.
Each choice is a nonempty residue class modulo a finite power of $p$.
The resulting common depth is $\max(d_0,e_0)$.
Only the smaller depth has been raised. There is no assertion that
both depths can be made arbitrarily large: $h$ and $h-1$ cannot both
be divisible by $p$. The selected contacts have $\mu=1$, and this
finite common depth is sufficient.

These new prime powers are coprime to every inherited annihilator
modulus and to one another. The Chinese remainder theorem therefore
preserves all original constraints. For the real return, one may first
enlarge the fixed progression modulus by the order of the component
group of the compact closure of the cyclic real orbit. After choosing
an actual permitted return as its residue, the closure of that fixed
progression is a coset of a connected compact torus. Multiplication by
any further positive integer is surjective on that torus. Every
subprogression whose residue preserves the fixed congruence consequently
has the same closure. The added prime-power constraints therefore
preserve the required real return. The fixed-prime neighborhoods are
preserved by their already retained congruences.
\end{proof}

The retained Gaussian annihilation congruences also put the required
half-generator numerator points in their respective formal groups.
In particular the $R_0$ annihilation required in~\cite{InverseVieta}
is retained together with the compatible moduli of~\cite[P16]{Uniform};
it is not inferred solely from divisibility of a group order by four.
The new prime-power congruences preserve these old annihilation
conditions. With the common depth $d\ge1$, all selected numerator depths equal
$d$, because their fixed nonzero root multipliers have been excluded
from the selected primes. Exact logarithmic labels and relative
order-four approximation now give
\[
 v(C_i)\ge4d,\qquad v(E_i)\ge2d,
 \qquad v(\mathcal G)\ge4d,\qquad v(\mathcal H)\ge2d,
 \qquad v(\mathcal N)\ge3d.
\]
The exact label Lucas residual vanishes. Every ordinary fourth-power
summand has valuation at least $8d-3>0$, every reciprocal-square
summand at least $4d-3>0$, and the input summand at least $24d-4>0$.
Thus every factor at the successful place is one modulo $p$.
At the opposite place its inverse is zero or one by the arbitrary-depth
separation proof. The actual reciprocal trace is therefore a unit of
residue one or two, and the prepared pole input has valuation one.

\section{Fixed places and completion of the witnesses}

The old ordered construction fixes compact neighborhoods on which all
point functions are finite, the reference reciprocal $f_H$ is nonzero,
and the scalar conditions persist. Shrinking these neighborhoods keeps
$E_i$, $\mathcal N$ and all homogeneous Lucas coordinates bounded.
The input value $\beta$ is fixed before these neighborhoods and the
comparison primes are chosen. The new progression refinement
preserves them. All powers of $q$ in the new factors are positive,
so at every fixed place $q\to0$ gives $W\to1$, $V\to2$ and
$b\to(204/409)^2$. The original local square classes and depths at
$2,3,17$ follow with sufficiently deep preparation: the inputs are
local squares with respective depths $4,2,2$, satisfying the retained
prepared-input requirements. At each prepared
real embedding, $q>0$ and all test powers are even, so all factors
are at least one and the actual trace is positive. The same limit
puts the actual input in $(0,17)$.

The unchanged contact factorization gives even valuations of $R(t_j)$
away from the selected primes. Division by the actual trace square preserves
parity. Thus the prepared pole inputs have square reduced denominators,
and the unchanged prepared pole-realization theorem applies.

At a selected prime the reciprocal trace is a unit, so the actual
pole input has valuation one, rather than a negative odd valuation.
At the other primes the even-parity argument applies. The two actual
inputs are now fixed and each receives its own prepared pole witnesses;
no earlier prime or reference choice is revisited. Fill the finite
point charts, the cross products, the input numerator, and the
homogeneous Lucas coordinates from these points. In particular
$P=L_3s_4$ gives $P s_3=2S s_4$. This supplies a rational solution
of the complete physical system.

The selected physical weights satisfy $q_0,q_1>0$. Hence $u=q_0$
and $v=q_1/q_0$ are nonzero and the additional unit coordinates
$v^{-3},u/v,u^{-1}/v$ have their required rational values. The
certified Laurent transformation therefore carries these physical
witnesses to the compiled quadratic system. All new prime-power
precisions, avoidance cutoffs and slopes may depend on the fixed
integer input and its previously chosen roots. They impose no new
input-dependent lower bound on $N$. One fixed threshold, combined
with the soundness threshold, proves Theorem~\ref{thm:arithmetic}.

\section{The defining quartic and its exact coefficients}\label{sec:quartic193-algebra}
Put
\[
 \begin{aligned}
 K_0&=126616122944180358034867798936,\\
 L_0&=484991541504995159662206970923766531374664106904986910720.
 \end{aligned}
\]
All variable counts include the rational input. The integer $C=149^N$ is
fixed and is not a quantified coordinate.

\begin{theorem}\label{thm:quartic193}
Under the arithmetic premises of Theorem~\ref{thm:arithmetic}, there is a
fixed integer $N_0$ such that, for every $N\ge N_0$, the polynomial
\[
 F_N\in\mathbb Z[X_1,\ldots,X_{193}]
\]
specified by the primitive quadratic table has rational input projection
exactly $\mathbb Z$. For every $N\ge1$, this polynomial is a sum of
167 primitive integral quadratic squares, has degree exactly four,
integer content one, exactly 2770 nonzero coordinate monomials, and
coefficient height
\[
                         H(F_N)=\max\{K_0,128\cdot149^{4N}\}.
\]
There is also a polynomial $\widetilde F_N$ with the same rational zero
set, the same number of variables and quadratic squares, degree four,
content one, exactly 2626 nonzero coordinate monomials, and height
\[
                  H(\widetilde F_N)=\max\{L_0,128\cdot149^{4N}\}.
\]
The coefficient crossovers are $N=4$ and $N=7$, respectively. Neither
crossover is asserted to be a sufficient defining threshold; $N_0$
remains ineffective.
\end{theorem}

The complete computation package is
\begin{quote}\small
\nolinkurl{computation/quartic/seven_units193_v74}.
\end{quote}
Its exact coefficient files are
\begin{quote}\small\raggedright
\nolinkurl{primitive/quadratic-family.json},\\
\nolinkurl{primitive/polynomial-family.json}.
\end{quote}
The alternative files are in \nolinkurl{fewer_monomials}.
Their coordinate order is the ascending
order of the retained source indices, with the input first. Each
coefficient record specifies its integer coefficient, power of the
fixed parameter $C$, and complete coordinate monomial. These records
specify every coefficient of both families.

\subsection{Physical storage and census}
The reference coordinates and scalar inverse Vieta point coordinates are
those of the retained source. For roles $1,3,4$ use the six-field
translated charts of \eqref{eq:chart}. Retain the primary scalar ratio
$A=L_2$ and the scalar quotient. The four alternate-ratio slots store
the cleared errors $E_i$ of \eqref{eq:cross}; their square slots store
$E_i^2$. The input uses $\mathcal N,\mathcal N^2,\mathcal N^4$.
The homogeneous Lucas circuit uses
five field coordinates $\mathcal C,\mathcal P,\mathcal Q,\mathcal H,
\mathcal H^2$. In particular the relation
\[
                         \mathcal P s_3=2Ss_4
\]
replaces both the stored primary ratio $L_3$ and its product with $s_4$.
The identities required for arithmetic extraction hold on every
consumed contact fiber; the proof does not require a global division by
$s_3$, $s_4$ or $\mathcal P$.

The coordinate census is
\begin{center}
\begin{tabular}{lr}
\hline
Part & Rational coordinates\\\hline
Input, rational normalizers, contact and pole coordinates &44\\
Contact units &7\\
Reference chart and reference outputs &10\\
Three translated numerator pairs &36\\
Scalar inverse Vieta point coordinates &14\\
Nonreference reciprocal-square outputs &8\\
Scalar primary ratio and scalar quotient &4\\
Four cleared cross comparisons and their squares &16\\
Fourth power of the first cleared cross comparison &2\\
Cleared input and its square and fourth power &6\\
Homogeneous Lucas circuit and its square &10\\
Factor evaluations and partial products &36\\\hline
Total &193\\\hline
\end{tabular}
\end{center}
Each field coordinate contributes two rational coordinates in the basis
$1,\sqrt2$. The coordinate-name table records the real and imaginary
storage conventions, including the half-imaginary square outputs.

\subsection{The unit chart and quadratic compilation}
For the exact coordinate transformation, write $u=q_0$ and
$v=q_1/q_0$. Before the final renumbering the seven units are
\[
\begin{array}{c|ccccccc}
\text{index}&77&90&91&92&901&902&903\\\hline
\text{value}&u&v&v^{-1}&u^{-1}v^{-2}&v^{-3}&uv^{-1}&u^{-1}v^{-1}.
\end{array}
\]
Retain the five literal quadratic equations
\begin{equation}\label{eq:193-unit-guards}
\begin{gathered}
 x_{77}x_{903}=x_{91},\qquad x_{90}x_{91}=1,\qquad
 x_{77}x_{92}=x_{91}^2,\\
 x_{90}x_{901}=x_{91}^2,\qquad
 x_{902}=x_{77}x_{91}.
\end{gathered}
\end{equation}
The equation $x_{90}x_{91}=1$ forces these two coordinates to be
nonzero. Then $x_{77}x_{903}=x_{91}$ forces $x_{77}$ and $x_{903}$
to be nonzero. The remaining three equations determine nonzero
coordinates $x_{92},x_{901},x_{902}$. Thus the seven-coordinate unit
chart is isomorphic over $\mathbb Q$ to the two-dimensional torus:
with $u=x_{77}$ and $v=x_{90}$ all seven values are exactly those
in the table. In particular the omitted coordinate $u^{-1}$ is
recoverable as $v x_{903}$ and introduces no boundary component.
This replaces the earlier eight-coordinate chart without changing
the physical equations or any arithmetic premise.

For each nonunit storage coordinate $z_i$, the certificate specifies an
integer pair $(a_i,b_i)$ and makes the substitution
\[
                         z_i=u^{a_i}v^{b_i}X_i.
\]
The input has pair $(0,0)$. Each equation is divided by its specified
unit monomial. If a term has nonunit monomial $m$, original unit
exponent $e$, and equation multiplier $r$, its remaining unit exponent is
\[
                    e+\sum_{i\in m}(a_i,b_i)-r.
\]
The exact certificate expresses this exponent as a sum of at most
$2-\deg m$ of the seven retained unit exponents. Replacing the unit
monomial by the corresponding coordinates therefore makes every term
an ordinary polynomial of degree at most two. The independent support
checker verifies all 804 literal and selected-test terms in the 122
rational or field equation groups. Splitting the field equations and
adjoining \eqref{eq:193-unit-guards} gives 167 rational quadrics.

The parameter factors are nonzero powers of $C$. Multiplying an equation
by such a factor and clearing its rational coefficient denominators
preserves its zero set. Each resulting integral row is divided by its
integer content. The file \nolinkurl{transformation.json} records every
coordinate weight, equation multiplier, parameter power, primitive
normalization and final coordinate number. Thus the compiled system and
the physical system used in the arithmetic proof have the same rational
input projection. The unit equations exclude every potential boundary
introduced by a Laurent denominator.

\subsection{The two reference bases}
The eight equations consisting of the reference-square definitions and
the six reference relations admit the constant basis change included in
\nolinkurl{source/reference-row-map.json}. Its determinant is
\[
             214450708242592116705963714180667354532074632082\ne0.
\]
All eight equations have the same Laurent row weight $(2,2)$.
Consequently the constant basis change commutes with their common
monomial division. Restoring the primitive reference basis gives
$F_N$; retaining the alternative basis gives $\widetilde F_N$.
The two full quadratic systems have the same rational zero set for every
fixed nonzero $C$. Summing their respective row squares therefore gives
the same rational zero set as well.

\subsection{Reducing the parameter degree}
The final parameter normalization is especially simple. In the
193-coordinate numbering put
\[
 \mathcal I=\{44,45,46,47,52,53,54,55,56,57,99,100\}.
\]
These coordinates store the square output for the scalar numerator,
the four squared cross products, and the fourth power of the first
cross product. Define an invertible diagonal map $D_C$ by
\[
 (D_C Z)_i=\begin{cases} C^{-1}Z_i,&i\in\mathcal I,\\
 Z_i,&i\notin\mathcal I.\end{cases}
\]
It fixes the input. For each primitive quadratic residual $q_j(C,X)$
define
\[
 \widehat q_j(C,Z)=C^{k_j}q_j(C,D_CZ),
\]
where $k_j$ is the negative of the smallest parameter exponent after
substitution. The explicit row table gives $k_j\in\{0,1,2\}$ and
all parameter exponents of $\widehat q_j$ in $\{0,1\}$. The same
substitution and row normalization apply to the alternative reference
basis. Let $\widehat F_N=\sum_j\widehat q_j(149^N,Z)^2$ and let
$\widehat{\widetilde F}_N$ denote the corresponding alternative sum.

\begin{theorem}\label{thm:compact193}
For every $N\ge1$, $\widehat F_N$ and $\widehat{\widetilde F}_N$
are integer quartics in 193 total coordinates, each a sum of 167
primitive integral quadratic squares. They have content one and
respectively 2770 and 2626 nonzero monomials. Their exact heights are
\[
 H(\widehat F_N)=\max\{K_0,149^{2N+2}\},\qquad
 H(\widehat{\widetilde F}_N)=\max\{L_0,149^{2N+2}\}.
\]
The coefficient crossovers are $N=6$ and $N=13$. Under the retained
arithmetic premises, both define $\mathbb Z$ for every sufficiently
large fixed $N$; no effective sufficient value is claimed.
\end{theorem}
\begin{proof}
The nonzero parameter makes $D_C$ invertible, and each row multiplier
$C^{k_j}$ is nonzero. Over $\mathbb Q$ a sum of squares vanishes
exactly when every summand vanishes. Consequently these operations
preserve the rational input projection on every fiber, including all
unit and point-chart boundaries already treated above.

The independent coefficient reconstruction checks every transformed
row and expands all 167 squares without importing the producer. The
resulting quartics have parameter degree two and respectively 2775
and 2631 parameter-separated terms. Every coordinate occurs. Their
input fourth-power coefficient is one, so their degree is four and
content is one. For every row the specialized coefficient gcd at
$C=149$ is one. This implies the same statement at every $149^N$:
for a prime other than 149, multiplying coefficients by powers of
149 cannot change simultaneous divisibility; divisibility by 149
is determined by the terms of parameter degree zero.

Here is a short exact height certificate. For either reference basis
write $J=K_0$ or $J=L_0$. The largest absolute constant coefficient
is $J$, and the largest quadratic coefficient is $22201=149^2$;
both pure coefficients occur. All pure linear coefficients have
absolute value at most
\[
 M=384444691473074,\qquad M^2\le22201J.
\]
There are exactly five coefficient polynomials with more than one
parameter power: $289C^2-578$, and four of the form $A+C^2$, with
\[
 A\in\{17086754668551163707,\ 204934086831195733699,\ 1,\ 1\}.
\]
This is the list for the primitive basis. In the alternative basis
the first value is instead
\[
 1747145497790223592810060466409331738555986526175401275,
\]
and the other three are unchanged. Each value satisfies
$22201A\le22200J$ for its corresponding $J$. For $C\ge149$, the exceptional
coefficient $289C^2-578$ is positive and at most $22201C^2$.
Writing $H=\max\{J,22201C^2\}$, we have
\[
 MC\le\sqrt{22201J}\,C\le H,
 \qquad
 A+C^2\le\frac{22200}{22201}J+
             \frac1{22201}(22201C^2)\le H.
\]
Thus every coefficient has absolute value at most $H$, and an
occurring pure coefficient attains it. This proves the exact height
for every real $C\ge149$, hence for $C=149^N$. The sole mixed-sign
coefficient is nonzero there, proving the stated constant support
counts. The crossover inequalities follow by exact integer comparison.
The rational equivalence gives the defining assertion from
Theorem~\ref{thm:arithmetic}.
\end{proof}

The improved coefficient tables are
\begin{quote}\small\raggedright
\nolinkurl{compact/quadratic-family.json},\\
\nolinkurl{compact/polynomial-family.json},\\
\nolinkurl{compact_fewer_monomials/quadratic-family.json},\\
\nolinkurl{compact_fewer_monomials/polynomial-family.json}.
\end{quote}
The two additional standard-library programs
\nolinkurl{build_scaled.py} and \nolinkurl{check_scaled.py}
produce and independently verify these tables. Their parameter
transformations record all twelve coordinate powers and every row
normalization.

\subsection{Expansion, height and independent verification}
An independent checker reconstructs all 81 new physical component rows
and all 81 retained source rows, verifies their exact monomial
pullbacks and the five unit guards, and expands the 167 row squares by
integer sparse arithmetic. The resulting expansions have respectively
2777 and 2633 parameter-separated terms, occupying 2770 and 2626
coordinate monomials. Every coordinate occurs.

Every coefficient polynomial has coefficients of one sign except
\[
                             289(C^2-2).
\]
This exception is nonzero for $C=149^N\ge149$, so the support counts are
valid for every $N\ge1$. The coefficient of $X_1^4$ is exactly one;
it proves both degree four and integer content one.

For the primitive basis, exact evaluation at $N=1,2,3$ gives height
$K_0$. Put $C_0=149^4$. For every coefficient polynomial $f(C)$, both
\[
                   128(C_0+Z)^4+f(C_0+Z),\qquad
                   128(C_0+Z)^4-f(C_0+Z)
\]
have nonnegative integer coefficients as polynomials in $Z$. These
5540 exact coefficient checks prove $|f(C)|\le128C^4$ for $C\ge C_0$.
A coefficient $128C^4$ occurs, so the bound is attained. The alternative
basis has the identical certificate with $C_0=149^7$, 5252 comparisons,
and exact evaluation at $N=1,\ldots,6$. A constant coefficient of absolute
value $K_0$ or $L_0$ occurs in the corresponding table. This proves both
height formulas and their precise crossovers.

The standalone commands are
\begin{quote}\ttfamily\raggedright
python build.py\\
python verify\_height.py\\
python build\_scaled.py\\
python check\_scaled.py\\
python check\_transfer.py\\
python check\_support.py\\
python check.py --folder primitive --primitive-reference\\
python check.py --folder fewer\_monomials
\end{quote}
The construction and support and height checks use only the Python
standard library. The independent physical-equation checker uses
SymPy 1.14.0 and imports neither the construction module nor an SMT
solver. The exact finite certificates prove the algebraic and coefficient
assertions. The rational defining property follows separately from
Theorem~\ref{thm:arithmetic} and its retained arithmetic premises.

\begin{thebibliography}{9}\raggedright
\bibitem{InverseVieta}
\emph{A quartic in 198 variables},
\nolinkurl{sources/h10_quartic198/arithmetic/inverse_vieta_198.tex},
especially the inverse-Vieta scalar chart, the unweighted height
aggregate, uniform scalar completeness and fresh-witness sections.

\bibitem{Uniform}
\emph{Existing 214: uniform existence and actual completeness interfaces},
\nolinkurl{sources/h10_quartic198/audit/source-uniform-existence.md},
P04, P07--P19, with the subsequent additive-aggregate and affine-label
changes specified in~\cite{InverseVieta,Quartic204}.

\bibitem{Quartic204}
\emph{A quartic with 204 variables},
\nolinkurl{sources/h10_quartic198/source/papers/quartic_existence/quartic204_20261114_v27.tex},
retained scalar and contact-height interfaces and half-integral
formal-group extraction.

\bibitem{Weighted}
\emph{Heights of positive weighted elliptic traces},
\nolinkurl{sources/h10_quartic198/source/papers/trace_heights/weighted_trace_20261114_v27.tex}.

\bibitem{Poles}
\emph{Congruent-twist poles},
\nolinkurl{sources/h10_quartic198/source/papers/pole_parameters/congruent_twist_poles_20261114_v27.tex},
with the retained prepared-realization interface of~\cite{Uniform}.

\bibitem{PrimeSeparation}
\emph{Separation and denominator conditions in quadratic prime patterns},
\nolinkurl{papers/prime_selection/prime_selection_20261109_v12.tex},
Theorem 2.2.
\end{thebibliography}

\end{document}
